% Generated draft: requires mathematical and source review.
\documentclass[11pt]{article}
\usepackage[margin=25mm]{geometry}
\usepackage{fontspec}
\setmainfont{Latin Modern Roman}
\usepackage{amsmath,amssymb,mathtools,unicode-math}
\setmathfont{Latin Modern Math}
\usepackage{enumitem,needspace,xurl,hyperref}
\setlength{\parindent}{0pt}
\setlength{\parskip}{5pt}
\setlength{\emergencystretch}{3em}
\newcommand{\N}{\mathbb{N}}
\newcommand{\Z}{\mathbb{Z}}
\newcommand{\Q}{\mathbb{Q}}
\newcommand{\R}{\mathbb{R}}
\newcommand{\C}{\mathbb{C}}
\newcommand{\F}{\mathbb{F}}
\newcommand{\A}{\mathbb{A}}
\newcommand{\PP}{\mathbb{P}}
\newcommand{\E}{\mathbb{E}}
\newcommand{\eps}{\varepsilon}
\newcommand{\dd}{\,\mathrm{d}}
\begin{document}
\section*{UnsolvedMath ID: 20003091}
\textbf{Source problem code:} catalogue entry
\subsection{Definitions and mathematical statement}
The source asks for an algorithm to find the maximally stretched lamination. Its extracted text also says that generically a Thurston geodesic between two Teichm\"uller points is not unique; this is retained as background rather than an algorithmic specification. The input encoding, precise lamination, output topology, and meaning of algorithm are not given.

For the supplied special-case discussion, let \(S=S_{1,1}\) be the once-punctured torus, let \(X,Y\) be distinct marked hyperbolic structures, and put
\[
 d_{\mathrm{Th}}(X,Y)=\log\sup_{\beta}\frac{\ell_Y(\beta)}{\ell_X(\beta)},
\]
where \(\beta\) ranges over essential simple closed curves. Let \(\Lambda(X,Y)\) denote the canonical largest maximally stretched chain-recurrent \emph{unmeasured} geodesic lamination in the convention of the supplied report. This object must be distinguished from a projective measured lamination maximizing the length ratio. Its complete intrinsic definition and Hausdorff-space conventions are not supplied.

The aid reports, for fixed \(X\), discontinuity of
\[
 Y\longmapsto\Lambda(X,Y)
\]
from numerical Fenchel--Nielsen coordinates to Hausdorff lamination output. At either boundary of a simple-curve out-envelope, an approach with output just a closed curve \(\alpha\) meets a boundary output consisting of \(\alpha\) and a coherently spiralling leaf. The two boundary types are denoted \(\alpha_0^+\) and \(\alpha_0^-\). The precise out-envelope and spiral conventions require the cited formulas.

Under the usual represented-space interpretation in which coordinate reals are supplied by Cauchy names and the desired output has the stated Hausdorff metric representation, a total Type-2 algorithm would induce a continuous function. The supplied discontinuity therefore obstructs such a total algorithm, already on this punctured-torus setting. This is not a claim that every symbolic representation, every weaker output, or every promised input subclass has the same obstruction.

The aid also reports boundary examples whose unique maximizing projective measured lamination is the same \([\alpha]\), but whose unmeasured outputs have different spiralling leaves. More explicitly, the set
\[
 \mathcal M(X,Y)=\left\{[\mu]:\mu\ne0\text{ measured},\ 
 \frac{\ell_Y(\mu)}{\ell_X(\mu)}=e^{d_{\mathrm{Th}}(X,Y)}\right\}
\]
is \(\{[\alpha]\}\) in both examples, while their \(\Lambda(X,Y)\) differ. Thus the full set of maximizing measured laminations alone does not recover the unmeasured spiral. This does not assert that complete marked metric data or the entire length-ratio function fail to determine it.

With stronger exact symbolic input, a certified unique simple-curve maximizer \(\alpha\), and decidable equality in the relevant exact comparisons, the aid attributes to DLRT's explicit twist formulas, equations (21)--(22), a terminating three-way classifier with output
\[
 \alpha,\qquad\alpha_0^+,\qquad\alpha_0^-.
\]
The symbolic input class, formulas, threshold comparisons, and certification procedure are not included in the record, so this is a reported promised-case classifier rather than an implementable general algorithm specified here. It does not decide equality from arbitrary numerical Cauchy names or settle every interpretation of the original algorithm question.
\subsection{Short English statement}
The supplied punctured-torus analysis separates numerical computability from exact symbolic classification. Hausdorff discontinuity rules out the stated total Cauchy-input algorithm, while a promised exact setting reportedly admits a three-way classifier. Maximizing measured laminations alone miss the unmeasured spiral data.
\subsection{Sources}
\begin{itemize}
\item Local UnsolvedMath registry, permanent ID \(20003091\), exact source code \texttt{catalogue entry}; canonical \texttt{aim-topology-notes.json}, \texttt{notes[178]}, tag \texttt{problem}.
\item AIM problem list ``Problems on Thurston metric,'' dated 21 April 2014, workshop ``Lipschitz metric on Teichm\"uller space,'' Problem 2.1: \url{https://aimath.org/pastworkshops/lipschitzteichproblems.pdf}. This URL is supplied as a citation only; no PDF was fetched. The record calls the statement exact, although its extracted text has a hyphenated word and trailing geodesic-background sentence.
\item The research aid cites DLRT and explicit twist formulas numbered (21)--(22), but does not expand the author abbreviation or supply a title, identifier, full citation, or the formulas. The original literature field is empty; an authoritative bibliographic reference for the special-case analysis remains missing.
\item Research attempt 1: result \texttt{partial\_result}, problem status at run \texttt{partially\_solved}, rigor label \texttt{proved}, novelty confidence \texttt{low}; classification \texttt{PARTIAL-PROGRESS}, registry status \texttt{partially\_solved}. Created and updated 14 August 2026. These are supplied assessments requiring review.
\item Attribution: Ulam AI / UnsolvedMath Contributors, CC BY 4.0. The machine-generated research aid is not peer reviewed and requires independent verification before citation.
\end{itemize}
\end{document}
